\subsubsection{Circuito puramente inductivo}

En esta experiencia se observa una coherencia entre los datos medidos y el circuito simulado. La corriente simulada a partir de la tensión medida se le corresponde un error menor al 5\% en comparación con la magnitud medida. Este error puede ser explicado por la consideración de idealidad de los instrumentos utilizados. 
Por otro lado, en la tabla (\textbf{REFERENCIAR TABLA CON LOS RESULTADOS DE Q,S,PHI, Z}) se presenta una potencia activa. Esto tiene sentido pues el inductor se modela como un inductor real, es decir que tiene una resistencia interna. Dicha resistencia interna se modela con una resistencia en serie de \SI{22.1}{\ohm}. Este razonamiento tambien explicapor qué el factor de potencia no es nulo. Sin embargo, se percibe una carga predominantemente inductiva, pues la potencia reactiva es positiva y es factor de potencia es bajo.

